\documentclass[11pt]{article}
\usepackage[margin=1in]{geometry}
\usepackage{amsmath,amssymb,mathtools}
\usepackage{booktabs}
\usepackage{enumitem}
\usepackage{siunitx}
\sisetup{group-separator={,},group-minimum-digits=3}

\title{\textbf{Quiz 1} \vspace{ 0.4cm }\\ECON 300: Labour Economics}
\author{Instructor: Muhebullah Karimzada}
\date{January 22,  2026}

\begin{document}
\maketitle

\section*{Multiple Choice (5 $\times$ 1 = 5 points)}

\textbf{Choose the best answer for each question.}

\begin{enumerate}
\item Which of the following individuals is classified as \textbf{unemployed} in official labour force statistics?

\begin{enumerate}
\item A full-time student not looking for work  
\item A retired worker  
\item A worker on temporary layoff expecting recall  
\item A discouraged worker who stopped searching  
\end{enumerate}

\item The labour force participation rate (LFPR) measures:

\begin{enumerate}
\item The fraction of the population that is employed  
\item The fraction of the working-age population that is employed or unemployed  
\item The fraction of workers who are unemployed  
\item The fraction of workers actively searching for jobs  
\end{enumerate}

\item In the labour--leisure model, an increase in the wage rate makes leisure:

\begin{enumerate}
\item Cheaper  
\item More valuable  
\item More expensive in terms of foregone income  
\item Irrelevant for labour supply  
\end{enumerate}

\item The substitution effect of a wage increase tends to:

\begin{enumerate}
\item Reduce hours worked  
\item Increase hours worked  
\item Leave hours unchanged  
\item Always dominate the income effect  
\end{enumerate}

\item Which of the following would most likely \textbf{increase} a worker's reservation wage?

\begin{enumerate}
\item A fall in non-labour income  
\item Lower commuting costs  
\item Stronger preference for leisure  
\item A higher market wage  
\end{enumerate}
\end{enumerate}



\subsection*{Problem 1: Labour Force Statistics and Interpretation (7.5 points)}

A city has a working-age population (15+) of 1{,}000{,}000 people. The breakdown is:

\begin{itemize}
\item Employed: 640{,}000  
\item Unemployed: 60{,}000  
\item Not in the labour force: 300{,}000  
\end{itemize}

\begin{enumerate}
\item[(a)] (2.5 points)  
Calculate the \textbf{labour force} and the \textbf{labour force participation rate}.

\item[(b)] (2.5 points)  
Calculate the \textbf{unemployment rate} and the \textbf{employment-to-population ratio}.

\item[(c)] (2.5 points)  
Suppose a recession causes:
\begin{itemize}
\item 20{,}000 employed workers to lose their jobs, and  
\item 15{,}000 discouraged workers to stop searching for work.
\end{itemize}

Recalculate the unemployment rate and explain why it may rise by \textbf{less than expected}.
\end{enumerate}

\subsection*{Problem 2: Labour Supply, Market Equilibrium, and Minimum Wage (7.5 points)}

A worker's weekly budget constraint is:

\[
C = 600 + 25h
\]

and utility is:

\[
U = C - 0.5h^2
\]

The market labour demand and supply are:

\[
L^D = 120 - 2w, 
\qquad 
L^S = 20 + 3w
\]

\begin{enumerate}
\item[(a)] (2.5 points)  
Find the worker's \textbf{optimal hours worked}.

\item[(b)] (2.5 points)  
Find the \textbf{equilibrium wage and employment} in the labour market.

\item[(c)] (2.5 points)  
A minimum wage of \( w = 20 \) is introduced.  
Determine whether it is binding and calculate the resulting \textbf{unemployment}.
\end{enumerate}

\vfill
\textbf{End of Quiz}

\end{document}
